\section{Implementation Details}

\paragraph{Appendix roadmap.}
The main paper compresses the evidence into the story that scale correction is directional and that the hard part is choosing the right scale regime. This appendix gives the details behind that claim. Section~\ref{sec:appendix_data} documents how ScaleQA-Natural is mined, audited, and separated from training and official evaluation. Section~\ref{sec:appendix_tables} reports development evidence showing that the scale-adapter gain is real but narrow. Section~\ref{sec:appendix_analysis} then follows the negative results: stronger prompts, verifiers, table calculators, candidate likelihood scoring, Bayesian fusion, and cross-backbone checks often expose the right answer in the candidate pool, but still fail to choose it reliably. This is why we frame adapters and verifiers as diagnostic probes rather than as a complete repair.

\paragraph{Adapters.}
The answer-tuned setting continues from the ChartQA/TextVQA answer-only LoRA adapter used in the evaluation pipeline. The answer-only adapters use rank $r=16$, LoRA $\alpha=32$, dropout 0.05, and target the language-model projection and MLP modules \texttt{q\_proj}, \texttt{k\_proj}, \texttt{v\_proj}, \texttt{o\_proj}, \texttt{gate\_proj}, \texttt{up\_proj}, and \texttt{down\_proj}. Answer-only training uses AdamW with learning rate $2\times 10^{-5}$. The final Qwen scale adapters use large-positive pairs with $K=300$, 300 optimization steps, AdamW learning rate $10^{-6}$, weight decay 0.01, gradient clipping at 1.0, $\beta=2.0$, and margin $m=0.10$. The clean direct setting starts from the raw Qwen2.5-VL-7B-Instruct backbone and applies only the scale-contrastive update. The InternVL2.5-1B adapter check uses the same pair construction, rank $r=8$, LoRA $\alpha=16$, dropout 0.05, 200 optimization steps, and the same language-model target modules.

\paragraph{Why separate answer-tuned and clean-direct settings.}
The answer-only adapter is not a scale intervention. It mostly makes the model better aligned to short chart answers, numeric formatting, and the evaluation prompt. Keeping this setting separate matters because the clean-direct scale-only adapter is unsafe on the raw backbone (Table~\ref{tab:clean_direct}), while the answer-tuned setting can absorb a small scale-specific update without the same global collapse. This difference suggests that scale correction is not a standalone numerical operation. It depends on a model already being in a chart-answering regime where answer extraction, formatting, and short-answer calibration are reasonably stable.

\paragraph{Decoding and scoring.}
All single-adapter systems use the same short-answer prompt and deterministic short-answer decoding template: the question is followed by ``Answer briefly.'' and generation uses greedy decoding with \texttt{do\_sample=False} and \texttt{max\_new\_tokens=24}. We evaluate numeric answers with the repository evaluator, which uses normalized exact matching for text and a 5\% numerical tolerance for numeric answers.

\paragraph{ScaleQA-Natural annotation.}
Candidate scale conflicts are mined from model predictions and automatically prefilled with numeric-ratio features. The review workflow assigns labels including scale diagnostic candidate, requires image inspection, non-scale, uncertain, and duplicate. Positive labels require that the gold answer and prediction refer to the same chart quantity but differ primarily by a simple scale relation; wrong entity, wrong operation, OCR failure, arbitrary numeric miss, ambiguous question, and duplicate cases are excluded or sent to adjudication. The final ScaleQA-Natural-520 set contains author-adjudicated positive candidates after question/gold deduplication. The image-evaluable subset contains 473 examples with available ChartQA image rows. The audit workflow exports three files: the 739-row candidate audit, a conflict queue, and a 108-row final-positive adjudication queue. After this adjudication, an independent second reviewer audited the 520 final positive inclusions and agreed with all inclusion labels. The release package is intended to include item IDs, source split and row index, question, gold answer, model prediction, subtype, direction, review label, audit metadata, and scripts for reproducing the filters. We reference ChartQA images by source row rather than redistributing them.

\paragraph{PlotQA external subset.}
The PlotQA external check uses a 500-example subset sampled from the official PlotQA test parquet after filtering to examples with a single numeric answer. We use the same relaxed numeric evaluator as the ChartQA experiments. This subset is intended as an external numeric-answer transfer check rather than a full PlotQA benchmark.

\section{Data Flow and Annotation Audit}
\label{sec:appendix_data}

ScaleQA-Natural is useful only if it is not mistaken for an ordinary benchmark split or for hidden training data. Table~\ref{tab:data_flow} therefore separates five roles: answer-only tuning, scale-pair training, development analysis, diagnostic set construction, and held-out or external evaluation. The reported Qwen scale adapters are trained from ChartQA training rows and training-time prediction caches; ScaleQA-Natural is used to diagnose regimes and evaluate policy selection, not to train those scale adapters. The held-out ChartQA test result remains the official aggregate task result.

\begin{table*}[h]
\centering
\small
\resizebox{\textwidth}{!}{\begin{tabular}{p{0.18\linewidth}p{0.25\linewidth}p{0.25\linewidth}p{0.24\linewidth}}
\toprule
Component & Source & Used for & Leakage control \\
\midrule
Answer-only tuning & ChartQA/TextVQA training data & Starting adapter for the main Qwen setting & Reported separately from scale tuning; not counted as a scale-objective gain \\
Scale-pair training & ChartQA training rows and training-set prediction/log-prob caches & LoRA scale adapters and pair-mode ablations & Official ChartQA test and PlotQA are not used for scale-adapter gradients \\
Development analysis & ChartQA validation and ScaleQA-Core diagnostics & Pair-mode selection, seed checks, and failure analysis & Reported as development evidence; held-out ChartQA test is the main official result \\
ScaleQA-Natural & Mined model outputs, machine audit, and author adjudication & Diagnostic regime analysis and router experiments & Used as a diagnostic set, not as ordinary task training data for the reported Qwen scale adapters \\
Held-out/external evaluation & ChartQA test, TextVQA-500, PlotQA numeric-500 & Final task result, safety, and external robustness check & Not used to train the reported scale adapters \\
\bottomrule
\end{tabular}
}
\caption{Data flow and split hygiene.}
\label{tab:data_flow}
\end{table*}

\begin{table}[h]
\centering
\small
\resizebox{\linewidth}{!}{\begin{tabular}{llrr}
\toprule
Audit view & Source / flag & Count & Share \\
\midrule
Image-evaluable 473 & train\_shard1 & 314 & 66.4\% \\
Image-evaluable 473 & val & 71 & 15.0\% \\
Image-evaluable 473 & train & 47 & 9.9\% \\
Image-evaluable 473 & test & 39 & 8.2\% \\
Image-evaluable 473 & train\_shard2 & 1 & 0.2\% \\
Image-evaluable 473 & train\_shard3 & 1 & 0.2\% \\
Final 520 & Contains ChartQA train occurrence & 336 & 64.6\% \\
Final 520 & Contains ChartQA validation occurrence & 59 & 11.3\% \\
Final 520 & Contains ChartQA test occurrence & 23 & 4.4\% \\
\bottomrule
\end{tabular}
}
\caption{ScaleQA-Natural source audit. The set is a diagnostic slice mined from multiple model outputs and source rows, not an ordinary training/test split or leaderboard. Occurrence flags are counted over the final 520 positives; source-split rows are counted over the 473 image-evaluable items.}
\label{tab:source_audit}
\end{table}

The annotation story has two parts. First, strict factor-of-100 failures are sparse in any single ChartQA output file (Table~\ref{tab:strict_prevalence}), so a single split underestimates the phenomenon. Second, mining introduces the opposite risk: plausible numeric ratios can be false positives when the model answered the wrong entity or operation. We therefore use ScaleQA-Natural as a diagnostic set whose positives require author adjudication and an independent inclusion audit (Table~\ref{tab:annotation_audit_app}), and we keep uncertain/rejected examples visible (Table~\ref{tab:negative_audit_examples}) rather than silently discarding them. The resulting set is deliberately conservative: it favors precision of the scale-conflict label over treating every ratio-like mismatch as a scale error.

\begin{table}[h]
\centering
\small
\resizebox{\linewidth}{!}{\begin{tabular}{lrl}
\toprule
Audit item & Count/Score & Notes \\
\midrule
Candidate pool & 739 & First-pass reviewed candidates \\
Observed agreement & 0.633 & Three-way labels \\
Cohen's $\kappa$ & 0.443 & Include/review/reject \\
Conflicts for review & 271 & Exported separately \\
Final positives & 520 & Author-adjudicated \\
Second-pass include & 412 & Deterministic audit \\
Second-pass review/reject & 108 & Manually adjudicated \\
Independent inclusion audit & 520 & Second human reviewer \\
Second-review agreement & 1.000 & Final positives; zero label changes \\
\bottomrule
\end{tabular}
}
\caption{Annotation audit for ScaleQA-Natural. The deterministic metadata pass flags conflicts for author adjudication; the final positive inclusions are then independently audited by a second human reviewer.}
\label{tab:annotation_audit_app}
\end{table}

\begin{table}[h]
\centering
\small
\resizebox{\linewidth}{!}{\begin{tabular}{llrr}
\toprule
Model / pipeline & Split & Strict errors & Rate \\
\midrule
Qwen2.5-VL-7B + answer/large & ChartQA test & 10 / 2500 & 0.40\% \\
Qwen2.5-VL-3B direct/large & ChartQA test & 19 / 2500 & 0.76\% \\
Qwen2.5-VL-3B direct/large & ChartQA val & 21 / 1920 & 1.09\% \\
InternVL2.5-1B & ChartQA test & 3 / 2500 & 0.12\% \\
InternVL2.5-1B & ChartQA val & 17 / 1920 & 0.89\% \\
\bottomrule
\end{tabular}
}
\caption{Strict factor-of-100 scale-error prevalence in available ChartQA validation/test prediction files. These sparse natural subsets motivate the larger ScaleQA-Natural diagnostic set.}
\label{tab:strict_prevalence}
\end{table}

\begin{table*}[h]
\centering
\small
\resizebox{\textwidth}{!}{\begin{tabular}{lll}
\toprule
Audit label & Auto subtype & Question snippet \\
\midrule
uncertain & percent\_to\_decimal & How much Good more than Bad? \\
uncertain & value\_scaled\_up & How many diamond cutting and polishing jobs were there in the Far East i... \\
uncertain & decimal\_to\_percent & How many respondents say no? \\
uncertain & decimal\_to\_percent & How many says that internet has good influence on education? \\
uncertain & decimal\_to\_percent & How confident are people about Angela Merkel doing the right thing regar... \\
uncertain & decimal\_to\_percent & How many more people answered Very acceptable for the most and least pop... \\
\bottomrule
\end{tabular}
}
\caption{Examples from the negative/uncertain audit queue. These candidates contain a plausible scale relation but are not included as positive diagnostics without stronger evidence that the prediction and gold answer refer to the same chart quantity.}
\label{tab:negative_audit_examples}
\end{table*}

\section{Additional Results}
\label{sec:appendix_tables}

The development results explain why the paper does not present the large-positive adapter as a general correction method. On ChartQA validation, the adapter recovers strict gold-scale errors much more often than it changes ordinary examples (Tables~\ref{tab:bootstrap} and~\ref{tab:seed_consistency}), but the aggregate confidence interval is small and overlaps zero on the full split. This is exactly the diagnostic pattern we care about: the intervention finds a real failure mode, but it is not a broad accuracy engine.

\begin{table*}[h]
\centering
\small
\resizebox{\textwidth}{!}{\begin{tabular}{lrrrrrrr}
\toprule
Split & N & Base & Ours & Delta & CI low & CI high & $p(\Delta \leq 0)$ \\
\midrule
All examples & 1920 & 0.860420 & 0.862500 & 0.002080 & -0.001250 & 0.005520 & 0.113500 \\
Numeric-answer examples & 1452 & 0.860880 & 0.862810 & 0.001930 & -0.002200 & 0.006060 & 0.187800 \\
Scale-keyword numeric examples & 1179 & 0.845630 & 0.847840 & 0.002210 & -0.002880 & 0.007460 & 0.202500 \\
Strict gold-scale errors & 23 & 0.000000 & 0.113040 & 0.113040 & 0.000000 & 0.252170 & 0.036800 \\
Non-gold-scale examples & 1897 & 0.870850 & 0.871590 & 0.000740 & -0.002210 & 0.003800 & 0.327100 \\
\bottomrule
\end{tabular}
}
\caption{Bootstrap analysis on ChartQA validation for one large-positive adapter snapshot. The full-set confidence interval overlaps zero, so we emphasize the directionality and diagnostic analyses rather than overclaiming aggregate accuracy.}
\label{tab:bootstrap}
\end{table*}

\begin{table}[h]
\centering
\small
\resizebox{\linewidth}{!}{\begin{tabular}{lrrrrr}
\toprule
Split & Base & Ours & Std & Delta & N \\
\midrule
All examples & 0.860400 & 0.864000 & 0.000600 & 0.003500 & 1920 \\
Strict gold-scale errors & 0.000000 & 0.208700 & 0.036400 & 0.208700 & 23 \\
Non-gold-scale examples & 0.870800 & 0.871900 & 0.000400 & 0.001100 & 1897 \\
\bottomrule
\end{tabular}
}
\caption{ChartQA validation subset analysis for the answer-tuned large-positive scale adapter. This was used as development evidence; the main official result is the held-out test result in Table~\ref{tab:test_subset_results}.}
\label{tab:val_subset_results}
\end{table}

\begin{table*}[h]
\centering
\small
\resizebox{\textwidth}{!}{\begin{tabular}{lrrrrrrrr}
\toprule
Split & N & Base & Mean fixes & Mean breaks & Any fix & All fix & Any break & All break \\
\midrule
All examples & 1920 & 0.860400 & 0.021400 & 0.010900 & 0.005700 & 0.003100 & 0.002600 & 0.001600 \\
Base-wrong examples & 268 & 0.000000 & 0.153000 & 0.000000 & 0.041000 & 0.022400 & 0.000000 & 0.000000 \\
Base-correct examples & 1652 & 1.000000 & 0.000000 & 0.012700 & 0.000000 & 0.000000 & 0.003000 & 0.001800 \\
Strict gold-scale errors & 23 & 0.000000 & 0.565200 & 0.000000 & 0.130400 & 0.087000 & 0.000000 & 0.000000 \\
Non-gold-scale examples & 1897 & 0.870800 & 0.014800 & 0.011100 & 0.004200 & 0.002100 & 0.002600 & 0.001600 \\
\bottomrule
\end{tabular}
}
\caption{Seed-level fix and break rates on ChartQA validation.}
\label{tab:seed_consistency}
\end{table*}

\begin{table}[h]
\centering
\small
\begin{tabular}{lrrr}
\toprule
Method & Pairs per kind & Accuracy & Delta \\
\midrule
Base & 0 & 0.850700 & 0.000000 \\
Large-25 & 25 & 0.856000 & 0.005300 \\
Large-50 & 50 & 0.856000 & 0.005300 \\
Large-100 & 100 & 0.855300 & 0.004700 \\
Large-200 & 200 & 0.856000 & 0.005300 \\
Large-300 & 300 & 0.855300 & 0.004700 \\
Large-400 & 400 & 0.855300 & 0.004700 \\
Large-600 & 600 & 0.856000 & 0.005300 \\
Large-800 & 800 & 0.856000 & 0.005300 \\
\bottomrule
\end{tabular}

\caption{Large-positive pair-size curve on the ChartQA-1500 setting. Gains appear with few pairs and then saturate.}
\label{tab:size_curve}
\end{table}

Table~\ref{tab:size_curve} is another reason we avoid a ``more LoRA pairs solves scale'' claim. Gains appear quickly and then saturate, while the subtype table in the main paper shows that the opposite direction requires a different policy. The ablation in Table~\ref{tab:ablation} should therefore be read as a probe-selection experiment: pair direction changes which mistakes are repaired and which safe cases are put at risk.

\begin{table*}[t]
\centering
\small
\resizebox{\textwidth}{!}{%
\begin{tabular}{lrrrrrr}
\toprule
Method & Overall & $\Delta$ & Gold-scale & $\Delta$ & Non-scale & $\Delta$ \\
\midrule
Base & 79.00 & +0.00 & 0.00 & +0.00 & 82.81 & +0.00 \\
Scale-aware prompt & 79.00 & +0.00 & 0.00 & +0.00 & 82.81 & +0.00 \\
Direct SFT (800) & 78.20 & -0.80 & 4.35 & +4.35 & 81.76 & -1.05 \\
Direct SFT (1500) & 77.40 & -1.60 & 4.35 & +4.35 & 80.92 & -1.89 \\
Mixed contrastive & 79.40 & +0.40 & 8.70 & +8.70 & 82.81 & +0.00 \\
Balanced contrastive & 79.20 & +0.20 & 4.35 & +4.35 & 82.81 & +0.00 \\
Large-only contrastive & 79.00 & +0.00 & 4.35 & +4.35 & 82.60 & -0.21 \\
Real contrastive & 79.40 & +0.40 & 8.70 & +8.70 & 82.81 & +0.00 \\
\textbf{Real contrastive (5 seeds)} & \textbf{79.52} & \textbf{+0.52} & 11.30 & +11.30 & 82.81 & +0.00 \\
\bottomrule
\end{tabular}
}
\caption{Ablation results on ChartQA-500. Gold-scale denotes strict 100$\times$ scale-error cases where the base model predicts an answer at the wrong numerical scale.}
\label{tab:ablation}
\end{table*}


\section{Additional Transfer and Analysis Tables}
\label{sec:appendix_analysis}

This section contains the main negative evidence behind the paper's conclusion. If scale correction were a simple numeric normalization problem, stronger backbones, larger candidate pools, chart-to-table calculators, or Bayesian policy fusion should close the gap. Instead, the same pattern recurs. The correct answer is often present under some candidate or policy oracle, but deployable selectors remain directionally biased. This supports the paper's central distinction: the phenomenon transfers more robustly than the intervention.

\begin{table}[h]
\centering
\small
\resizebox{\linewidth}{!}{\begin{tabular}{llrrr}
\toprule
Backbone & Split & N & Accuracy & Strict scale errors \\
\midrule
Qwen2.5-VL-7B & ChartQA test & 2500 & 0.7804 & 10 \\
InternVL2.5-1B & ChartQA validation & 1920 & 0.7745 & 17 \\
InternVL2.5-1B & ChartQA test & 2500 & 0.7524 & 3 \\
InternVL2.5-1B & ScaleQA-Natural-520 & 473 & 0.3340 & -- \\
\bottomrule
\end{tabular}
}
\caption{Cross-backbone strict scale-like error diagnostic.}
\label{tab:cross_backbone}
\end{table}

\begin{table*}[h]
\centering
\small
\resizebox{\textwidth}{!}{\begin{tabular}{llrrrr}
\toprule
Set & Model & Base & Large & Small & Bal. \\
\midrule
ScaleQA-Natural-520 & InternVL2.5-1B & 0.3340 & 0.3446 & 0.3446 & -- \\
ScaleQA-Natural-520 & InternVL2.5-8B & 0.4672 & 0.4685 & 0.4723 & 0.4689 \\
ScaleQA-Natural-520 & InternVL2.5-26B & 0.4989 & -- & -- & -- \\
ScaleQA-stress1000 & InternVL2.5-8B & 0.7320 & 0.7272 & 0.7276 & 0.7256 \\
ScaleQA-stress1000 & Qwen2.5-VL-7B & 0.6940 & 0.6940 & 0.6940 & 0.6940 \\
PlotQA-557 & InternVL2.5-8B & 0.6912 & 0.6829 & 0.6862 & 0.6840 \\
PlotQA-500 & InternVL2.5-26B & 0.7200 & -- & -- & -- \\
ChartBench-500 & Qwen2.5-VL-7B & 0.4820 & 0.4820 & 0.4820 & 0.4820 \\
ChartBench-500 & InternVL2.5-8B & 0.3520 & 0.3512 & 0.3516 & 0.3520 \\
ChartBench-500 & InternVL2.5-26B & 0.4080 & -- & -- & -- \\
MathVista-243 & Qwen2.5-VL-7B & 0.6708 & 0.6708 & 0.6708 & 0.6708 \\
MathVista-243 & InternVL2.5-8B & 0.6337 & 0.6379 & 0.6387 & 0.6387 \\
MathVista numeric & InternVL2.5-26B & 0.6749 & -- & -- & -- \\
InfographicVQA-300 & Qwen2.5-VL-7B & 0.7767 & 0.7767 & 0.7767 & 0.7767 \\
InfographicVQA-300 & InternVL2.5-8B & 0.6067 & 0.6087 & 0.6073 & 0.6093 \\
InfographicVQA-300 & InternVL2.5-26B & 0.6300 & -- & -- & -- \\
\bottomrule
\end{tabular}
}
\caption{Full cross-backbone and external numeric transfer table.}
\label{tab:transfer_boundary_full}
\end{table*}

Tables~\ref{tab:cross_backbone}--\ref{tab:external_internvl8b} show the transfer boundary. Larger InternVL models improve the base level on ScaleQA-Natural and external numeric subsets, but the same large/small/balanced scale adapters do not become uniformly helpful. On ScaleQA-stress1000, for example, InternVL2.5-8B adapters slightly reduce accuracy relative to the base model; on MathVista and InfographicVQA the changes are small; on PlotQA the adapter can hurt. These results do not refute the existence of scale inconsistency. They show that a global scale intervention learned in one model/task setting is not the right abstraction for deployment.

\begin{table}[h]
\centering
\small
\resizebox{\linewidth}{!}{\begin{tabular}{lrrr}
\toprule
Backbone and setting & N & Accuracy & Std. \\
\midrule
InternVL2.5-1B base & 473 & 0.3340 & -- \\
InternVL2.5-1B large-positive LoRA & 473 & 0.3446 & -- \\
InternVL2.5-1B small-positive LoRA & 473 & 0.3446 & -- \\
\midrule
InternVL2.5-8B base & 473 & 0.4672 & -- \\
InternVL2.5-8B large-positive LoRA & 473 & 0.4685 & 0.0017 \\
InternVL2.5-8B small-positive LoRA & 473 & 0.4723 & 0.0035 \\
InternVL2.5-8B balanced LoRA & 473 & 0.4689 & 0.0025 \\
\bottomrule
\end{tabular}
}
\caption{InternVL scale adaptation on ScaleQA-Natural-520.}
\label{tab:internvl_adapter}
\end{table}

\begin{table}[h]
\centering
\small
\resizebox{\linewidth}{!}{\begin{tabular}{lrrr}
\toprule
InternVL2.5-8B setting & N & Accuracy & Std. \\
\midrule
Base & 1000 & 0.7320 & -- \\
Large-positive adapter & 1000 & 0.7272 & 0.0016 \\
Small-positive adapter & 1000 & 0.7276 & 0.0034 \\
Balanced adapter & 1000 & 0.7256 & 0.0023 \\
\bottomrule
\end{tabular}
}
\caption{InternVL2.5-8B on a synthetic ScaleQA-stress1000 set.}
\label{tab:internvl8b_stress}
\end{table}

\begin{table}[h]
\centering
\small
\begin{tabular}{lr}
\toprule
Prompting strategy & Accuracy \\
\midrule
Direct answer & 0.5243 \\
Axis-aware prompt & 0.5307 \\
Math-check prompt & 0.5137 \\
Scale prompt & 0.4905 \\
Self-check prompt & 0.5455 \\
Unit-guard prompt & 0.5180 \\
Majority vote & 0.5264 \\
Oracle prompt choice & 0.6744 \\
\bottomrule
\end{tabular}

\caption{Training-free prompting baselines on ScaleQA-Natural-520.}
\label{tab:tta_baselines_app}
\end{table}

\begin{table}[h]
\centering
\small
\resizebox{\linewidth}{!}{\begin{tabular}{lccc}
\toprule
Method & Overall & Pred. too large & Pred. too small \\
\midrule
Candidate majority & 0.5159 & 0.3029 & 0.7371 \\
Free-form verifier & 0.4588 & 0.1411 & 0.7888 \\
Letter verifier & 0.4651 & 0.1411 & 0.8017 \\
Scale-guard letter verifier & 0.4609 & 0.1452 & 0.7888 \\
Candidate oracle & 0.8288 & 0.7593 & 0.9009 \\
\bottomrule
\end{tabular}
}
\caption{Multi-candidate verifier on ScaleQA-Natural-520. Candidate oracle selects the best answer from the same candidate pool and is an upper bound.}
\label{tab:candidate_verifier}
\end{table}

The selector and verifier results give the sharpest version of the bottleneck. Self-check prompting is already a strong fixed policy, and the best no-leak selector improves on it only modestly (Table~\ref{tab:deployable_selectors}). Yet the candidate and policy oracles are much higher. The gap is not caused by a failure to generate alternatives: it is caused by failures to rank alternatives under chart-specific scale evidence. The structured verifier makes this especially clear. It is strong when the prediction is too small but nearly collapses when the prediction is too large (Table~\ref{tab:structured_verifier_direction}), even though the augmented oracle is high in both directions.

\begin{table*}[h]
\centering
\small
\resizebox{\textwidth}{!}{\begin{tabular}{lcc}
\toprule
Baseline or selector & Accuracy & Notes \\
\midrule
Self-check prompt & 0.5514 & Best fixed prompt on held-out splits \\
v8 no-leak selector & 0.5432 $\pm$ 0.0501 & Question, answer, numeric, and chart-context features \\
v9 expanded selector & 0.5501 $\pm$ 0.0460 & Adds structured/pairwise policies but does not improve \\
v10b no-leak selector & 0.5667 $\pm$ 0.0561 & Best clean selector; adds corrected TTA-v2 evidence \\
v11 log-prob selector & 0.5667 $\pm$ 0.0561 & Log-prob evidence preserves but does not improve v10b \\
v13 rationale selector & 0.5667 $\pm$ 0.0561 & Rationale-verifier evidence preserves v10b \\
v14/v15 candidate-likelihood selectors & 0.5667 $\pm$ 0.0561 & Answer/large/base likelihood evidence preserves v10b \\
v16/v17 top-8 selectors & 0.5667 $\pm$ 0.0561 & Larger candidate pools preserve v10b but raise only oracle \\
Statistical policy fusion & 0.5309 $\pm$ 0.0470 & Beta/Wilson reliability fusion \\
Bucketed Bayesian selector & 0.5239 $\pm$ 0.0480 & Bucketed policy priors; over-switches \\
Bayesian-MLP selector & 0.5218 $\pm$ 0.0540 & Bucket priors plus hidden-layer selector; over-switches \\
Pairwise top-8 verifier & 0.5433 & Candidate comparison prompt \\
Candidate likelihood top-8 & 0.4503 & Log-probability candidate scoring \\
Structured-context verifier & 0.4841 & Image, chart context, and top-4 candidates \\
Structured scale-audit verifier & 0.4376 & Explicit unit/scale audit prompt \\
Table extraction + calculator & 0.4228 & DePlot-style chart-to-table evidence, then symbolic calculation \\
\midrule
Original policy-pool oracle & 0.7023 $\pm$ 0.0384 & Upper bound over early policy pool \\
Expanded policy-pool oracle & 0.9004 & Upper bound after adding prompts/verifiers \\
Candidate top-8 oracle & 0.9175 & Correct answer often appears in candidate set \\
Augmented structured oracle & 0.9260 & Upper bound after adding scaled candidate variants \\
\bottomrule
\end{tabular}
}
\caption{Deployable selectors and strong baselines on ScaleQA-Natural. Learned and rule-based selectors use held-out splits, while fixed prompting, verifier, and table-calculator rows are full-set scores. The gap to the oracle rows shows that policy selection remains the bottleneck.}
\label{tab:deployable_selectors}
\end{table*}

\begin{table}[h]
\centering
\small
\resizebox{\linewidth}{!}{\begin{tabular}{lccc}
\toprule
Quantity & Mean & Std. & Additional statistic \\
\midrule
No-leak selector & 0.5667 & 0.0561 & 50 held-out splits \\
Self-check prompt & 0.5514 & 0.0430 & same splits \\
Selector - self-check & 0.0154 & 0.0492 & 95\% CI [0.0017, 0.0286] \\
Best fixed policy & 0.5714 & 0.0362 & split oracle over fixed policies \\
Policy oracle & 0.9004 & 0.0238 & upper bound \\
Switch rate & 48.9\% & 39.3\% & not using self-check \\
\bottomrule
\end{tabular}
}
\caption{No-leak selector stability on ScaleQA-Natural held-out splits. The improvement over self-check is modest but positive under a split-level bootstrap; the oracle gap remains much larger.}
\label{tab:selector_stability}
\end{table}

\begin{table*}[h]
\centering
\small
\resizebox{\textwidth}{!}{\begin{tabular}{lrrrrrr}
\toprule
Group & N & Selector & Ans. & Large & Self-check & Best fixed \\
\midrule
Direction: prediction too large & 241 & 0.4860 & 0.2946 & 0.1452 & 0.4357 & 0.4357 \\
Direction: prediction too small & 232 & 0.6534 & 0.7328 & 0.8621 & 0.6595 & 0.8621 \\
Source: test & 39 & 0.4275 & 0.3590 & 0.5385 & 0.3846 & 0.5385 \\
Source: val & 71 & 0.6250 & 0.5915 & 0.6338 & 0.6479 & 0.6479 \\
Source: train & 47 & 0.6303 & 0.6170 & 0.5319 & 0.6809 & 0.6809 \\
Source: train shard1 & 314 & 0.5587 & 0.4904 & 0.4522 & 0.5191 & 0.5191 \\
Subtype: decimal to percent & 109 & 0.5193 & 0.2752 & 0.1468 & 0.3853 & 0.3853 \\
Subtype: percent to decimal & 74 & 0.6411 & 0.8243 & 0.9324 & 0.7297 & 0.9324 \\
Subtype: value scaled down & 118 & 0.6578 & 0.7034 & 0.8729 & 0.6102 & 0.8729 \\
Subtype: value scaled up & 102 & 0.4098 & 0.2549 & 0.0784 & 0.4216 & 0.4216 \\
\bottomrule
\end{tabular}
}
\caption{Selector and fixed-policy behavior by source, direction, and subtype. ``Selector'' is the v10b held-out selector score averaged over test-fold appearances. Fixed-policy columns are item-level diagnostics from the candidate pool and are included to show directional reversal, not as a new benchmark leaderboard.}
\label{tab:selector_group_diagnostics}
\end{table*}

\begin{table}[h]
\centering
\small
\resizebox{\linewidth}{!}{\begin{tabular}{lccc}
\toprule
Method & Overall & Pred. too large & Pred. too small \\
\midrule
Structured-context verifier & 0.4503 & 0.0954 & 0.8190 \\
Structured scale-audit verifier & 0.4376 & 0.0996 & 0.7888 \\
Augmented structured oracle & 0.9260 & 0.9419 & 0.9095 \\
\bottomrule
\end{tabular}
}
\caption{Structured verifier by scale direction. Adding chart context and scaled candidate variants makes the oracle high, but deployable verifiers remain directionally biased.}
\label{tab:structured_verifier_direction}
\end{table}

\begin{table}[h]
\centering
\small
\resizebox{\linewidth}{!}{\begin{tabular}{lrrr}
\toprule
Subset & Reason & Calculator & Either \\
\midrule
All examples & 0.3890 & 0.4080 & 0.4228 \\
Percent-to-decimal & 0.8378 & 0.8243 & 0.8378 \\
Value scaled down & 0.6949 & 0.8220 & 0.8475 \\
Decimal-to-percent & 0.0275 & 0.0183 & 0.0183 \\
Value scaled up & 0.0098 & 0.0098 & 0.0098 \\
\bottomrule
\end{tabular}
}
\caption{DePlot-style table-extraction plus calculator baseline by subtype on ScaleQA-Natural. The method is strong for some regimes but fails on the opposite direction, again showing that scale correction is regime-dependent.}
\label{tab:table_calculator_subtypes}
\end{table}

The table-calculator baseline is a useful stress test because it is closer to symbolic reasoning than to answer-style prompting. It performs well on percent-to-decimal and value-scaled-down cases, but almost never solves decimal-to-percent or value-scaled-up cases (Table~\ref{tab:table_calculator_subtypes}). That asymmetry mirrors the adapter story: changing the evidence source does not remove the need to identify the regime first.

\begin{table}[h]
\centering
\small
\resizebox{\linewidth}{!}{\begin{tabular}{lccccc}
\toprule
Policy pool & Best fixed & Learned & +Chart ctx. & Subtype oracle & Inst. oracle \\
\midrule
Prompt regimes (473) & 0.5408 & 0.5366 & 0.5358 & 0.6192 & 0.6701 \\
Adapters+prompts (403) & 0.5032 & 0.4467 & 0.4615 & 0.8010 & 0.8410 \\
\bottomrule
\end{tabular}
}
\caption{Learned policy routing with question/base-answer numeric features and optional chart-side context extracted from the image. Values are 50-seed means.}
\label{tab:chart_context_router}
\end{table}

\begin{table}[h]
\centering
\small
\resizebox{\linewidth}{!}{\begin{tabular}{llr}
\toprule
Backbone & Method & Accuracy \\
\midrule
Qwen2.5-VL-7B & Answer-tuned & 0.902 \\
Qwen2.5-VL-7B & Answer-tuned + large-positive scale & 0.898 \\
Qwen2.5-VL-7B & Base + scale prompt & 0.832 \\
Qwen2.5-VL-7B & TTA majority & 0.892 \\
Qwen2.5-VL-7B & TTA oracle prompt choice & 0.904 \\
InternVL2.5-1B & Base & 0.600 \\
\bottomrule
\end{tabular}
}
\caption{External numeric PlotQA-500 check.}
\label{tab:plotqa_external}
\end{table}

\begin{table}[h]
\centering
\small
\resizebox{\linewidth}{!}{\begin{tabular}{lrrrrr}
\toprule
Dataset & N & Base & Large & Small & Balanced \\
\midrule
PlotQA & 557 & 0.6912 & 0.6829 & 0.6862 & 0.6840 \\
ChartBench & 500 & 0.3520 & 0.3512 & 0.3516 & 0.3520 \\
MathVista & 243 & 0.6337 & 0.6379 & 0.6387 & 0.6387 \\
InfographicVQA & 300 & 0.6067 & 0.6087 & 0.6073 & 0.6093 \\
\bottomrule
\end{tabular}
}
\caption{External numeric checks with InternVL2.5-8B.}
\label{tab:external_internvl8b}
\end{table}

\begin{table}[h]
\centering
\small
\resizebox{\linewidth}{!}{\begin{tabular}{lrrrrl}
\toprule
Subtype & N & Base & Large-pos. & Small-pos. & Best direction \\
\midrule
Decimal to percent & 109 & 0.0917 & 0.1560 & 0.7890 & Small \\
Value scaled up & 102 & 0.0784 & 0.0784 & 0.7059 & Small \\
Percent to decimal & 74 & 0.9054 & 0.9324 & 0.0270 & Large \\
Value scaled down & 118 & 0.8475 & 0.8559 & 0.0593 & Large \\
\bottomrule
\end{tabular}
}
\caption{ScaleQA-Natural subtype analysis on the 403-example reviewed set.}
\label{tab:subtype_results}
\end{table}

The subtype and taxonomy tables explain why the task is hard for deployable selectors. Most image-evaluable examples contain explicit percent evidence, but the correct action is not always ``add a percent'' or ``multiply by 100''. In decimal-to-percent and value-scaled-up cases, the small-positive policy is often best; in percent-to-decimal and value-scaled-down cases, the large-positive policy is best (Table~\ref{tab:subtype_results}). The same words---percent, average, difference, value lookup---appear across regimes (Table~\ref{tab:fine_taxonomy}), so a selector needs to reason over the chart, the question, the base answer, and the candidate answer jointly rather than keying off a single unit token.

\begin{table}[h]
\centering
\small
\resizebox{\linewidth}{!}{\begin{tabular}{lrr}
\toprule
Group & Count & Share \\
\midrule
\multicolumn{3}{l}{\textit{Question pattern}} \\
Percent / percentage & 133 & 28.1 \\
Average / mean & 75 & 15.9 \\
Value lookup & 75 & 15.9 \\
Difference / comparison & 65 & 13.7 \\
Maximum / highest & 44 & 9.3 \\
Minimum / lowest & 30 & 6.3 \\
\midrule
\multicolumn{3}{l}{\textit{Unit evidence}} \\
Percent in axis or question & 439 & 92.8 \\
No explicit unit & 14 & 3.0 \\
Large-unit word & 13 & 2.7 \\
Ratio or rate & 6 & 1.3 \\
Currency & 1 & 0.2 \\
\bottomrule
\end{tabular}
}
\caption{Fine-grained taxonomy of the image-evaluable ScaleQA-Natural subset. Counts summarize question patterns and unit evidence; categories are not mutually exclusive across the two blocks.}
\label{tab:fine_taxonomy}
\end{table}

\begin{table}[h]
\centering
\small
\resizebox{\linewidth}{!}{\begin{tabular}{lrrrrrr}
\toprule
Evaluation set & N & Base & Direct-large & Delta & Direct-small & Delta \\
\midrule
ChartQA validation & 1920 & 0.7536 & 0.7219 & -0.0318 & 0.6948 & -0.0589 \\
ChartQA test & 2500 & 0.7804 & 0.7456 & -0.0348 & 0.7136 & -0.0668 \\
ScaleQA-Core & 403 & 0.4591 & 0.4367 & -0.0223 & 0.4243 & -0.0347 \\
InternVL-mined & 192 & 0.1510 & 0.0938 & -0.0573 & 0.3542 & +0.2031 \\
\bottomrule
\end{tabular}
}
\caption{Clean direct scale-only adaptation from the raw Qwen2.5-VL-7B backbone.}
\label{tab:clean_direct}
\end{table}

Finally, Table~\ref{tab:clean_direct} is the safety counterexample. A scale-only update from the raw Qwen backbone degrades ChartQA validation, ChartQA test, and ScaleQA-Core, even though it helps a separate InternVL-mined subset in the opposite direction. This is why the paper treats scale adapters as probes of directional asymmetry, not as a universal repair. The case studies in Table~\ref{tab:case_study} give the qualitative version of the same lesson: the model can repair 0.61$\rightarrow$61 and still over-correct a true small decimal such as 0.0593.

\begin{table*}[h]
\centering
\small
\begin{tabular}{lp{0.47\linewidth}lll}
\toprule
Type & Question & Gold & Base & Calibrated \\
\midrule
Recovered & How much times Madagascar greater than Central Europe data? & 32 & 0.032 & 32 \\
Recovered & What the difference in value between Asia and Caribbean? & 1.18 & 0.18 & 1.18 \\
Recovered & What is the highest value of gray bar? & 61 & 0.61 & 61 \\
Recovered & What is the difference in \% between Poland and Australia for people who are worried about losing their job? & 17.05 & 0.1705 & 17.05 \\
Over-correction & Find out the average of the bottom three countries ?? & 0.0593 & 0.06 & 6.03333 \\
\bottomrule
\end{tabular}

\caption{Representative recoveries and an over-correction failure.}
\label{tab:case_study}
\end{table*}

\paragraph{Overall appendix takeaway.}
Taken together, the appendix strengthens the main claim in a deliberately conservative way. Scale-like mistakes are not frequent enough in a single benchmark split to support sweeping accuracy claims, but they are systematic enough to reveal a reliability failure once mined and audited. Global adapters expose the directionality, training-free self-checks provide a strong practical baseline, and expanded candidate pools show that many failures are recoverable in principle. The unresolved problem is selecting the right regime without gold subtype labels. This is the reason the paper argues for regime-aware scale reliability rather than presenting unconditional scale correction as solved.
